\begin{lemma}
\label{lemma-lift-pseudo-coherent}
Let $X$ be an affine scheme and let $U \subset X$ be a quasi-compact
open subscheme. For any pseudo-coherent object $E$ of $D(\mathcal{O}_U)$
there exists a bounded above complex of finite free $\mathcal{O}_X$-modules 
whose restriction to $U$ is isomorphic to $E$.
\end{lemma}

\begin{proof}
By Lemma \ref{lemma-pseudo-coherent} we see that $E$ is an object of
$D_\QCoh(\mathcal{O}_U)$. By
Lemma \ref{lemma-lift-quasi-coherent}
we may assume $E = E'|U$ for some object $E'$ of
$D_\QCoh(\mathcal{O}_X)$.
Write $X = \Spec(A)$. By Lemma \ref{lemma-affine-compare-bounded}
we can find a complex $M^\bullet$ of $A$-modules whose associated
complex of $\mathcal{O}_X$-modules is a representative of $E'$.

\medskip\noindent
Choose $f_1, \ldots, f_r \in A$ such that $U = D(f_1) \cup \ldots \cup D(f_r)$.
By Lemma \ref{lemma-pseudo-coherent-affine} the complexes
$M^\bullet_{f_j}$ are pseudo-coherent complexes of $A_{f_j}$-modules.
Let $n$ be an integer. Assume we have a map of complexes
$\alpha : F^\bullet \to M^\bullet$ where $F^\bullet$ is
bounded above, $F^i = 0$ for $i < n$, each $F^i$ is a finite free
$R$-module, such that
$$
H^i(\alpha_{f_j}) : H^i(F^\bullet_{f_j}) \to H^i(M^\bullet_{f_j})
$$
is an isomorphism for $i > n$ and surjective for $i = n$. Picture
$$
\xymatrix{
& F^n \ar[r] \ar[d]^\alpha & F^{n + 1} \ar[d]^\alpha \ar[r] & \ldots \\
M^{n-1} \ar[r] & M^n \ar[r] & M^{n + 1} \ar[r] & \ldots
}
$$
Since each $M^\bullet_{f_j}$ has vanishing cohomology
in large degrees we can find such a map for $n \gg 0$.
By induction on $n$ we are going to extend this to a map
of complexes $F^\bullet \to M^\bullet$
such that $H^i(\alpha_{f_j})$ is an isomorphism
for all $i$. The lemma will follow by taking $\widetilde{F^\bullet}$.

\medskip\noindent
The induction step will be to extend the diagram
above by adding $F^{n - 1}$. Let $C^\bullet$ be the cone on $\alpha$
(Derived Categories, Definition \ref{derived-definition-cone}).
The long exact sequence of cohomology shows that
$H^i(C^\bullet_{f_j}) = 0$ for $i \geq n$. By
More on Algebra, Lemma \ref{more-algebra-lemma-cone-pseudo-coherent}
we see that $C^\bullet_{f_j}$ is $(n - 1)$-pseudo-coherent. By
More on Algebra, Lemma \ref{more-algebra-lemma-finite-cohomology}
we see that $H^{n - 1}(C^\bullet_{f_j})$ is a finite $A_{f_j}$-module.
Choose a finite free $A$-module $F^{n - 1}$ and an $A$-module
$\beta : F^{n - 1} \to C^{n - 1}$ such that the composition
$F^{n - 1} \to C^{n - 1} \to C^n$ is zero and such that
$F^{n - 1}_{f_j}$ surjects onto $H^{n - 1}(C^\bullet_{f_j})$.
(Some details omitted; hint: clear denominators.)
Since $C^{n - 1} = M^{n - 1} \oplus F^n$
we can write $\beta = (\alpha^{n - 1}, -d^{n - 1})$. The vanishing of the
composition $F^{n - 1} \to C^{n - 1} \to C^n$ implies
these maps fit into a morphism of complexes
$$
\xymatrix{
& F^{n - 1} \ar[d]^{\alpha^{n - 1}} \ar[r]_{d^{n - 1}} &
F^n \ar[r] \ar[d]^\alpha &
F^{n + 1} \ar[d]^\alpha \ar[r] & \ldots \\
\ldots \ar[r] &
M^{n - 1} \ar[r] & M^n \ar[r] & M^{n + 1} \ar[r] & \ldots
}
$$
Moreover, these maps define a morphism of distinguished triangles
$$
\xymatrix{
(F^n \to \ldots) \ar[r] \ar[d] &
(F^{n-1} \to \ldots) \ar[r] \ar[d] &
F^{n-1} \ar[r] \ar[d]_\beta &
(F^n \to \ldots)[1] \ar[d] \\
(F^n \to \ldots) \ar[r] &
M^\bullet \ar[r] &
C^\bullet \ar[r] &
(F^n \to \ldots)[1]
}
$$
Hence our choice of $\beta$ implies that the map of complexes
$(F^{-1} \to \ldots) \to M^\bullet$ induces an isomorphism on
cohomology localized at $f_j$ in degrees $\geq n$ and a surjection
in degree $n - 1$. This finishes the proof of the lemma.
\end{proof}
